% Generated draft: requires mathematical and source review.
\documentclass[11pt]{article}
\usepackage[margin=25mm]{geometry}
\usepackage{fontspec}
\setmainfont{Latin Modern Roman}
\usepackage{amsmath,amssymb,mathtools,unicode-math}
\setmathfont{Latin Modern Math}
\usepackage{enumitem,needspace,xurl,hyperref}
\setlength{\parindent}{0pt}
\setlength{\parskip}{5pt}
\setlength{\emergencystretch}{3em}
\newcommand{\N}{\mathbb{N}}
\newcommand{\Z}{\mathbb{Z}}
\newcommand{\Q}{\mathbb{Q}}
\newcommand{\R}{\mathbb{R}}
\newcommand{\C}{\mathbb{C}}
\newcommand{\F}{\mathbb{F}}
\newcommand{\A}{\mathbb{A}}
\newcommand{\PP}{\mathbb{P}}
\newcommand{\E}{\mathbb{E}}
\newcommand{\eps}{\varepsilon}
\newcommand{\dd}{\,\mathrm{d}}
\begin{document}
\section*{UnsolvedMath ID: 30005905}
\textbf{Source problem code:} catalogue entry
\subsection{Definitions and mathematical statement}
A differential graded (dg) category over a field \(k\) is a category whose morphism spaces are cochain complexes of \(k\)-vector spaces, with composition compatible with the differentials. The source does not specify the base field or the admissibility hypotheses on its dg categories.

Write \((\mathcal D,\widetilde F)\) schematically for a dg categorical dynamical system, and let \((\mathcal T,F)\) denote its corresponding categorical dynamical system, where \(F\) is the exact endofunctor induced by the dg quasi-endofunctor \(\widetilde F\). The construction of \(\mathcal T\) from \(\mathcal D\) is not specified in the supplied excerpt. Denote the source's Hochschild homology entropy and Hochschild cohomology entropy by \(h_{HH_*}(\widetilde F)\) and \(h_{HH^*}(\widetilde F)\), respectively, and its categorical entropy by \(h_{\mathrm{cat}}(F)\). These symbols name the quantities in the question; their defining growth formulas, finiteness assumptions, and normalization conventions are absent from the supplied record.

\textbf{Problem.} Determine hypotheses on a dg categorical dynamical system under which its Hochschild homology entropy, its Hochschild cohomology entropy, or both coincide with the entropy of the corresponding categorical dynamical system. Schematically, characterize the systems for which
\[
 h_{HH_*}(\widetilde F)=h_{\mathrm{cat}}(F),
 \qquad\text{or}\qquad
 h_{HH^*}(\widetilde F)=h_{\mathrm{cat}}(F),
\]
with both equalities requested if ``(co)homology'' is intended simultaneously. This is a request for conditions ensuring equality, not an assertion of equality for every system. The excerpt does not determine whether \(h_{\mathrm{cat}}\) means the value \(h_0\) of a parameterized categorical entropy \(h_t\); no equality for all \(t\) is asserted here. A complete mathematical formulation requires the omitted entropy definitions and domain hypotheses.

\textbf{Supplied contextual qualifications.} The excerpt attributes higher representation-finite algebras to Iyama and Oppermann, higher representation-infinite and higher hereditary algebras to Herschend, Iyama and Oppermann, and twisted fractionally Calabi--Yau algebras to Herschend and Iyama. It reports that higher hereditary algebras are either higher representation-finite or higher representation-infinite, and that higher representation-finite algebras are examples of twisted fractionally Calabi--Yau algebras. These are background claims attributed to the excerpt's references [6]--[8], whose bibliographic entries are not supplied.

The excerpt reports calculations for the Serre functor \(S\) and its dg quasi-functor lift \(\widetilde S\). Under its twisted fractionally Calabi--Yau hypothesis, it gives
\[
 h_t(S)=rt,\qquad h_t^{\mathrm{pol}}(S)=0,
\]
and, with the additional assumption that the algebra is elementary,
\[
 h_{HH_*}(\widetilde S)=h_{HH^*}(\widetilde S)=0.
\]
Here \(r\) stands for the fractional Calabi--Yau parameter: the extracted text prints both the parameter and coefficient as ``pq'', apparently representing \(p/q\), but the typography and parameter conventions require checking. These formulas are preserved as supplied contextual claims, not independently verified theorems. In particular, the additional elementary-algebra hypothesis must be retained for the stated Hochschild entropy conclusion. The excerpt does not provide the definition of elementary or the full ambient hypotheses of this theorem.

The excerpt also says that, at the time of the source discussion, it was unknown whether every higher representation-finite algebra, or more generally every twisted fractionally Calabi--Yau algebra, is fractionally Calabi--Yau; it points to [1, Question 1.6]. This is a historical source qualification, not an independently verified current status.
\subsection{Short English statement}
Determine when the growth measured by Hochschild homology or cohomology agrees with categorical entropy for a dg dynamical system. The supplied record does not specify enough definitions or hypotheses to turn this into a unique fully quantified conjecture. It reports special Serre-functor results and labels the general question partially solved; those claims require human mathematical and source review.
\subsection{Sources}
\begin{itemize}
\item Local supplied registry record: \texttt{data/unsolvedmath/problems.json}, permanent UnsolvedMath ID \(30005905\), source problem code \texttt{catalogue entry}, title ``Hochschild and Categorical Entropy in dg Dynamics.'' No individual proposer is identified in the record.
\item Source report: \emph{Hochschild (Co)Homology and Applications} (2024), Oberwolfach Reports, pp. 1097--1162, DOI 10.4171/owr/2024/20. \url{https://doi.org/10.4171/owr/2024/20}. The supplied excerpt attributes the entropy question to [9, Question 2.13]; the bibliographic identity of [9] is not included.
\item Supplied literature reference: Han, ``entropies of Serre functors.'' \url{https://doi.org/10.1016/j.jalgebra.2024.03.031}. This is the attribution and descriptive label supplied by the registry.
\item Supplied literature reference: Kikuta and Ouchi, ``Hochschild entropy.'' \url{https://doi.org/10.1007/s40598-022-00210-5}. This is the attribution and descriptive label supplied by the registry.
\item Literature assessment dated exactly \texttt{2026-08-22}: the supplied curation labels the problem partially solved and states that equality of Hochschild and categorical entropy is proved for Serre functors in broad higher-hereditary and twisted fractionally Calabi--Yau cases, while the general dg-categorical question remains open. This is supplied dated curation derived from the OWR record, not an independently verified literature assessment. Its precise hypotheses and time-sensitive claims require checking.
\item The excerpt additionally cites [1, Question 1.6], [2, 2.6.1], [3, Remark 6.3], [6, Theorem 1.1], [7, Theorem 3.4], and [8]. Their complete bibliographic entries are absent; no identities or external links are inferred here.
\end{itemize}
\end{document}
